\section{Movimiento relativo}

Cuando observamos el movimiento de un objeto, normalmente describimos
su posición, velocidad y aceleración respecto de algún sistema de
referencia.

Sin embargo, un mismo objeto puede parecer que se mueve de maneras
diferentes cuando es observado desde sistemas de referencia distintos.

Por ejemplo, imaginemos a una persona que camina por el pasillo de un
tren en movimiento.

Para otro pasajero sentado dentro del tren, la persona puede avanzar
lentamente hacia la parte delantera.

En cambio, para alguien que permanece en una estación, esa misma
persona se desplaza con una velocidad diferente, porque además de
caminar participa del movimiento del tren.

También podemos considerar dos automóviles que circulan por una
carretera. Aunque ambos estén en movimiento respecto del suelo,
uno puede parecer prácticamente en reposo cuando se observa desde
el otro.

Estas situaciones muestran una característica fundamental de la
cinemática:

\begin{idea}[El movimiento depende del sistema de referencia]

No basta con afirmar que un objeto tiene una determinada velocidad.

Para describir correctamente su movimiento debemos indicar
\textbf{respecto de qué sistema de referencia} se mide esa velocidad.

Un mismo objeto puede tener velocidades diferentes respecto de
observadores distintos, sin que exista contradicción física.

\end{idea}

El objetivo de este capítulo es aprender a relacionar las posiciones,
velocidades y aceleraciones medidas desde diferentes sistemas de
referencia.


\section{El carácter relativo del movimiento}

En los capítulos anteriores establecimos sistemas de referencia para
describir movimientos rectilíneos, movimientos en varias dimensiones,
proyectiles y movimientos circulares.

En muchos ejemplos elegimos como referencia el suelo y lo
consideramos aproximadamente fijo.

Pero esta elección no es obligatoria.

Podemos estudiar un movimiento respecto de:

\begin{itemize}
    \item el suelo;
    \item un automóvil;
    \item un tren;
    \item una embarcación;
    \item una aeronave;
    \item otro objeto que también se encuentre en movimiento.
\end{itemize}

La posición y la velocidad asignadas al objeto pueden ser diferentes
en cada uno de estos sistemas.

\begin{definicion}[Movimiento relativo]

El movimiento relativo es la descripción del cambio de posición de
un cuerpo respecto de otro cuerpo o sistema de referencia.

\end{definicion}

Esta definición no significa que cualquier descripción sea
arbitraria.

Una vez elegido el sistema de referencia, las posiciones,
velocidades y aceleraciones deben determinarse mediante relaciones
matemáticas coherentes.

\begin{interpretacion}

Decir que un automóvil se mueve a
\(\SI{20}{\meter\per\second}\) tiene un significado incompleto
si no especificamos la referencia.

Habitualmente se sobreentiende que esa velocidad se mide respecto
del suelo.

Pero también podríamos medir la velocidad del automóvil respecto
de otro vehículo, y obtener un resultado diferente.

\end{interpretacion}


\section{Observadores y sistemas de referencia}

En Física utilizamos la palabra \textbf{observador} para referirnos
a quien realiza o describe mediciones desde un determinado sistema
de referencia.

No es necesario imaginar a una persona observando directamente
el movimiento.

Un sistema de referencia puede estar asociado a instrumentos
de medición, cámaras o sistemas de navegación.

Para estudiar el movimiento relativo necesitamos, al menos,
dos sistemas de referencia.

Consideremos:

\begin{itemize}
    \item \(S\): sistema de referencia fijo al suelo;
    \item \(S'\): sistema de referencia asociado a un vehículo.
\end{itemize}

El sistema \(S'\) puede moverse respecto del sistema \(S\).

Un mismo objeto \(P\) tendrá una posición medida desde \(S\)
y otra posición medida desde \(S'\).

\begin{modelo}[Sistemas de referencia en traslación]

Inicialmente estudiaremos sistemas cuyos ejes cartesianos
permanecen paralelos entre sí.

El origen de uno puede trasladarse respecto del otro, pero
no consideraremos rotaciones relativas de los ejes.

Además, utilizaremos la descripción clásica o galileana
del movimiento, apropiada cuando las velocidades involucradas
son mucho menores que la velocidad de la luz.

\end{modelo}

Esta elección permitirá construir las relaciones del movimiento
relativo sin introducir todavía efectos asociados a sistemas
de referencia en rotación.


\section{Notación para las magnitudes relativas}

Cuando intervienen varios observadores es indispensable
establecer una notación que evite ambigüedades.

Utilizaremos:

\[
\vect{r}_{A/B}
\]

para representar la posición de \(A\) respecto de \(B\).

Asimismo:

\[
\vect{v}_{A/B}
\]

representará la velocidad de \(A\) respecto de \(B\).

Finalmente:

\[
\vect{a}_{A/B}
\]

representará la aceleración de \(A\) respecto de \(B\).

\begin{recordatorio}[Lectura de subíndices]

La expresión:

\[
\vect{v}_{A/B}
\]

se lee:

\begin{center}
``velocidad de \(A\) respecto de \(B\)''.
\end{center}

El primer subíndice identifica el objeto cuyo movimiento
se describe.

El segundo identifica la referencia.

\end{recordatorio}

No debemos confundir:

\[
\vect{v}_{A/B}
\]

con:

\[
\vect{v}_{B/A}.
\]

Son magnitudes relacionadas, pero sus sentidos son opuestos
cuando ambas se expresan mediante ejes paralelos.


\section{Posición relativa entre dos partículas}

Comenzaremos por el caso más sencillo.

Supongamos que dos partículas, \(A\) y \(B\), ocupan
diferentes posiciones en un mismo sistema cartesiano.

Sus vectores de posición son:

\[
\vect{r}_A
\]

y:

\[
\vect{r}_B.
\]

Queremos determinar dónde se encuentra \(A\) cuando tomamos
como referencia la posición de \(B\).

\begin{center}
\begin{tikzpicture}[scale=0.95]

    \coordinate (O) at (0,0);
    \coordinate (A) at (5.3,3.2);
    \coordinate (B) at (2.1,1.1);

    \draw[-{Latex},thick]
        (-0.5,0) -- (6.6,0)
        node[right] {\(x\)};

    \draw[-{Latex},thick]
        (0,-0.5) -- (0,4.4)
        node[above] {\(y\)};

    \draw[-{Latex},very thick,AzulPrincipal]
        (O) -- (A)
        node[midway,above left] {\(\vect{r}_A\)};

    \draw[-{Latex},very thick,VerdePrincipal]
        (O) -- (B)
        node[midway,below right] {\(\vect{r}_B\)};

    \draw[-{Latex},very thick,RojoPrincipal]
        (B) -- (A)
        node[midway,below right] {\(\vect{r}_{A/B}\)};

    \fill (O) circle (2pt);
    \fill (A) circle (2.5pt);
    \fill (B) circle (2.5pt);

    \node[below left] at (O) {\(O\)};
    \node[above right] at (A) {\(A\)};
    \node[above left] at (B) {\(B\)};

\end{tikzpicture}
\end{center}

Geométricamente:

\[
\vect{r}_B+\vect{r}_{A/B}=\vect{r}_A.
\]

Despejando:

\[
\boxed{
\vect{r}_{A/B}=\vect{r}_A-\vect{r}_B
}.
\]

Esta es la relación fundamental de posición relativa entre
dos partículas medidas desde un mismo sistema de referencia.

\begin{interpretacion}

El vector:

\[
\vect{r}_{A/B}
\]

apunta desde \(B\) hacia \(A\).

No representa necesariamente el movimiento realizado
por ninguna de las partículas.

Representa su separación vectorial en un instante determinado.

\end{interpretacion}


\section{Posición relativa mediante componentes rectangulares}

En dos dimensiones:

\[
\vect{r}_A=x_A\ihat+y_A\jhat,
\]

\[
\vect{r}_B=x_B\ihat+y_B\jhat.
\]

Entonces:

\[
\vect{r}_{A/B}
=
(x_A-x_B)\ihat+(y_A-y_B)\jhat.
\]

Por tanto:

\[
\boxed{
x_{A/B}=x_A-x_B
},
\]

\[
\boxed{
y_{A/B}=y_A-y_B
}.
\]

En tres dimensiones:

\[
\boxed{
\vect{r}_{A/B}
=
(x_A-x_B)\ihat
+
(y_A-y_B)\jhat
+
(z_A-z_B)\khat
}.
\]

La magnitud de la separación es:

\[
\boxed{
r_{A/B}
=
\sqrt{
(x_A-x_B)^2+
(y_A-y_B)^2+
(z_A-z_B)^2
}
}.
\]

En dos dimensiones se omite el término correspondiente a \(z\).

\begin{advertencia}

La separación entre dos partículas es una magnitud escalar
no negativa.

El vector posición relativa, en cambio, contiene dirección
y sentido.

Por ello no debemos reemplazar automáticamente:

\[
\vect{r}_{A/B}
\]

por:

\[
r_{A/B}.
\]

\end{advertencia}


\section{Ejemplo: posición relativa en dos dimensiones}

\begin{ejemplo}[Localización de dos vehículos]

En un sistema cartesiano fijo al suelo, dos vehículos
se encuentran en las posiciones:

\[
\vect{r}_A
=
(120\ihat+80\jhat)\,\si{\meter},
\]

\[
\vect{r}_B
=
(40\ihat+140\jhat)\,\si{\meter}.
\]

Determine:

\begin{enumerate}[label=\alph*)]
    \item La posición de \(A\) respecto de \(B\).
    \item La distancia entre ambos vehículos.
    \item La dirección de \(A\) respecto de \(B\).
\end{enumerate}

\end{ejemplo}

\subsection*{Posición relativa}

Partimos de:

\[
\vect{r}_{A/B}=\vect{r}_A-\vect{r}_B.
\]

Por componentes:

\[
x_{A/B}=120-40=\SI{80}{\meter},
\]

\[
y_{A/B}=80-140=-\SI{60}{\meter}.
\]

Por tanto:

\[
\boxed{
\vect{r}_{A/B}
=
(80\ihat-60\jhat)\,\si{\meter}
}.
\]

\subsection*{Distancia}

La magnitud es:

\[
r_{A/B}
=
\sqrt{80^2+(-60)^2}\,\si{\meter}.
\]

Entonces:

\[
\boxed{
r_{A/B}=\SI{100}{\meter}
}.
\]

\subsection*{Dirección}

El vector tiene:

\[
x_{A/B}>0,
\qquad
y_{A/B}<0.
\]

Por tanto, apunta hacia el cuarto cuadrante.

La magnitud del ángulo respecto de \(+x\) es:

\[
\theta
=
\arctan\left(\frac{60}{80}\right).
\]

Así:

\[
\boxed{
\theta\approx36.9\grado
}
\]

por debajo del eje \(+x\).

\begin{interpretacion}

Desde la posición de \(B\), el vehículo \(A\) está
\(\SI{80}{\meter}\) hacia \(+x\) y
\(\SI{60}{\meter}\) hacia \(-y\).

La distancia recta entre ellos es \(\SI{100}{\meter}\).

\end{interpretacion}


\section{Velocidad relativa entre dos partículas}

La posición relativa puede cambiar con el tiempo.

Si ambas partículas se mueven, tenemos:

\[
\vect{r}_{A/B}(t)
=
\vect{r}_A(t)-\vect{r}_B(t).
\]

La velocidad relativa es la tasa instantánea de cambio
de la posición relativa:

\[
\vect{v}_{A/B}
=
\deriv{\vect{r}_{A/B}}{t}.
\]

Derivando:

\[
\vect{v}_{A/B}
=
\deriv{\vect{r}_A}{t}
-
\deriv{\vect{r}_B}{t}.
\]

Como:

\[
\vect{v}_A=\deriv{\vect{r}_A}{t},
\]

\[
\vect{v}_B=\deriv{\vect{r}_B}{t},
\]

obtenemos:

\[
\boxed{
\vect{v}_{A/B}
=
\vect{v}_A-\vect{v}_B
}.
\]

\begin{definicion}[Velocidad relativa]

La velocidad relativa de \(A\) respecto de \(B\) es
la tasa de cambio temporal de la posición de \(A\)
respecto de \(B\).

Para partículas medidas desde un mismo sistema de
referencia cartesiano:

\[
\boxed{
\vect{v}_{A/B}=\vect{v}_A-\vect{v}_B
}.
\]

\end{definicion}

Su unidad SI es:

\[
\si{\meter\per\second}.
\]

\begin{interpretacion}

La velocidad relativa describe cómo cambia la separación
vectorial entre dos partículas.

No indica necesariamente la rapidez con la que cambia
la distancia escalar entre ellas.

Esta diferencia resulta importante en movimientos
bidimensionales.

\end{interpretacion}


\section{Componentes rectangulares de la velocidad relativa}

Si:

\[
\vect{v}_A=v_{Ax}\ihat+v_{Ay}\jhat,
\]

\[
\vect{v}_B=v_{Bx}\ihat+v_{By}\jhat,
\]

entonces:

\[
\vect{v}_{A/B}
=
(v_{Ax}-v_{Bx})\ihat
+
(v_{Ay}-v_{By})\jhat.
\]

Por tanto:

\[
\boxed{
v_{A/B,x}=v_{Ax}-v_{Bx}
},
\]

\[
\boxed{
v_{A/B,y}=v_{Ay}-v_{By}
}.
\]

En tres dimensiones también:

\[
\boxed{
v_{A/B,z}=v_{Az}-v_{Bz}
}.
\]

La magnitud es:

\[
\boxed{
v_{A/B}
=
\sqrt{
v_{A/B,x}^2+
v_{A/B,y}^2+
v_{A/B,z}^2
}
}.
\]

Esta es una relación vectorial y debe resolverse por componentes
cuando el movimiento ocurre en varias dimensiones.


\section{Cambio del orden de los observadores}

Para estudiar la velocidad de \(B\) respecto de \(A\):

\[
\vect{v}_{B/A}=\vect{v}_B-\vect{v}_A.
\]

Comparando:

\[
\vect{v}_{A/B}=\vect{v}_A-\vect{v}_B,
\]

obtenemos:

\[
\boxed{
\vect{v}_{B/A}=-\vect{v}_{A/B}
}.
\]

Por tanto:

\[
\boxed{
v_{B/A}=v_{A/B}
}.
\]

Las magnitudes de ambas velocidades relativas son iguales,
pero sus sentidos son opuestos.

\begin{interpretacion}

Si desde el automóvil \(B\) observamos que el automóvil \(A\)
se mueve hacia la derecha a cierta rapidez relativa, desde
\(A\) observaremos que \(B\) se mueve hacia la izquierda
con esa misma rapidez relativa, siempre que comparemos
mediciones expresadas sobre ejes paralelos.

\end{interpretacion}


\section{Movimiento relativo en una dimensión}

Comenzaremos con movimientos sobre una misma línea recta.

Elegimos:

\[
+x
\]

hacia la derecha.

Las velocidades pueden tener valores positivos o negativos.

La relación es:

\[
\boxed{
v_{A/B,x}=v_{Ax}-v_{Bx}
}.
\]

Esta única ecuación permite analizar movimientos:

\begin{itemize}
    \item en el mismo sentido;
    \item en sentidos contrarios;
    \item con uno de los objetos en reposo.
\end{itemize}

No es necesario memorizar reglas independientes para cada situación.


\section{Objetos que se mueven en el mismo sentido}

Consideremos dos automóviles que se desplazan hacia \(+x\).

Sus velocidades respecto del suelo son:

\[
v_A=\SI{25}{\meter\per\second},
\]

\[
v_B=\SI{18}{\meter\per\second}.
\]

La velocidad de \(A\) respecto de \(B\) es:

\[
v_{A/B}=v_A-v_B.
\]

Entonces:

\[
v_{A/B}=25-18.
\]

Por tanto:

\[
\boxed{
v_{A/B}=+\SI{7}{\meter\per\second}
}.
\]

\begin{interpretacion}

Para una persona que se desplaza junto con \(B\), el
automóvil \(A\) avanza hacia \(+x\) a
\(\SI{7}{\meter\per\second}\).

Aunque ambos vehículos se mueven rápidamente respecto
del suelo, la velocidad relativa es menor.

\end{interpretacion}


\section{Objetos que se mueven en sentidos contrarios}

Consideremos ahora:

\[
v_A=+\SI{20}{\meter\per\second},
\]

\[
v_B=-\SI{15}{\meter\per\second}.
\]

La velocidad relativa es:

\[
v_{A/B}=v_A-v_B.
\]

Sustituyendo:

\[
v_{A/B}=20-(-15).
\]

Entonces:

\[
\boxed{
v_{A/B}=+\SI{35}{\meter\per\second}
}.
\]

\begin{interpretacion}

En este caso las magnitudes parecen sumarse porque las
velocidades tienen signos opuestos.

Pero no hemos utilizado una regla especial de suma.

Hemos aplicado exactamente la misma relación general:

\[
v_{A/B}=v_A-v_B.
\]

\end{interpretacion}


\section{Velocidad relativa y rapidez de acercamiento}

Cuando dos objetos se mueven sobre la misma línea, suele
interesarnos saber si se están acercando o alejando.

Esta pregunta requiere analizar no sólo la velocidad
relativa, sino también la posición relativa.

Definimos:

\[
x_{A/B}=x_A-x_B.
\]

La velocidad relativa es:

\[
v_{A/B}=\deriv{x_{A/B}}{t}.
\]

La distancia entre los objetos es:

\[
d=|x_{A/B}|.
\]

Cuando:

\[
x_{A/B}\neq0,
\]

la tasa de cambio de la distancia es:

\[
\boxed{
\deriv{d}{t}
=
\frac{x_{A/B}}{|x_{A/B}|}
v_{A/B}
}.
\]

Si:

\[
\deriv{d}{t}<0,
\]

la distancia disminuye y los objetos se acercan.

Si:

\[
\deriv{d}{t}>0,
\]

la distancia aumenta y se alejan.

\begin{advertencia}

La rapidez relativa no debe confundirse automáticamente
con la rapidez de acercamiento.

Incluso en una dimensión es necesario saber qué objeto
se encuentra delante.

En dos dimensiones, la distinción es todavía más importante.

\end{advertencia}


\section{Ejemplo: persecución entre automóviles}

\begin{ejemplo}[Encuentro en una carretera]

Dos automóviles se mueven sobre una carretera recta
en la dirección \(+x\).

En \(t=0\), el automóvil \(A\) se encuentra en:

\[
x_{A0}=\SI{0}{\meter},
\]

con velocidad:

\[
v_A=\SI{25}{\meter\per\second}.
\]

El automóvil \(B\) está inicialmente en:

\[
x_{B0}=\SI{150}{\meter},
\]

con velocidad:

\[
v_B=\SI{15}{\meter\per\second}.
\]

Ambos mantienen velocidades constantes.

Determine cuánto tiempo tarda \(A\) en alcanzar a \(B\)
y dónde ocurre el encuentro.

\end{ejemplo}

\subsection*{Modelo físico}

Utilizamos movimiento rectilíneo uniforme para ambos vehículos.

El eje \(+x\) apunta en el sentido de sus movimientos.

Las posiciones son:

\[
x_A(t)=x_{A0}+v_At,
\]

\[
x_B(t)=x_{B0}+v_Bt.
\]

Por tanto:

\[
x_A(t)=25t,
\]

\[
x_B(t)=150+15t.
\]

\subsection*{Condición de encuentro}

Los automóviles se encuentran cuando:

\[
x_A=x_B.
\]

Entonces:

\[
25t=150+15t.
\]

Despejando:

\[
10t=150.
\]

Por tanto:

\[
\boxed{
t=\SI{15.0}{\second}
}.
\]

\subsection*{Posición del encuentro}

Sustituimos en la ecuación de \(A\):

\[
x_A=25(15.0).
\]

Entonces:

\[
\boxed{
x=\SI{375}{\meter}
}.
\]

Comprobamos con \(B\):

\[
x_B=150+15(15.0)=\SI{375}{\meter}.
\]

Ambas posiciones coinciden.

\begin{interpretacion}

El automóvil \(A\) se acerca a \(B\) con rapidez relativa:

\[
25-15=\SI{10}{\meter\per\second}.
\]

Como la separación inicial es \(\SI{150}{\meter}\),
necesita \(\SI{15}{\second}\) para eliminarla.

El método relativo proporciona el mismo resultado que
el análisis de ambas posiciones.

\end{interpretacion}


\section{Aceleración relativa entre dos partículas}

La velocidad relativa también puede cambiar con el tiempo.

Partimos de:

\[
\vect{v}_{A/B}=\vect{v}_A-\vect{v}_B.
\]

Derivamos:

\[
\deriv{\vect{v}_{A/B}}{t}
=
\deriv{\vect{v}_A}{t}
-
\deriv{\vect{v}_B}{t}.
\]

Por tanto:

\[
\boxed{
\vect{a}_{A/B}=\vect{a}_A-\vect{a}_B
}.
\]

\begin{definicion}[Aceleración relativa]

La aceleración relativa de \(A\) respecto de \(B\) es
la tasa de cambio temporal de la velocidad relativa.

Para partículas descritas mediante un mismo sistema
cartesiano:

\[
\boxed{
\vect{a}_{A/B}=\vect{a}_A-\vect{a}_B
}.
\]

\end{definicion}

En componentes:

\[
\boxed{
a_{A/B,x}=a_{Ax}-a_{Bx}
},
\]

\[
\boxed{
a_{A/B,y}=a_{Ay}-a_{By}
},
\]

\[
\boxed{
a_{A/B,z}=a_{Az}-a_{Bz}
}.
\]

La unidad SI es:

\[
\si{\meter\per\second\squared}.
\]


\section{Aceleración relativa nula}

Si dos partículas tienen la misma aceleración:

\[
\vect{a}_A=\vect{a}_B,
\]

entonces:

\[
\boxed{
\vect{a}_{A/B}=\vect{0}
}.
\]

Esto significa que la velocidad relativa permanece constante.

Sin embargo, no implica que ambas partículas tengan
la misma velocidad.

Por ejemplo, dos cuerpos que caen libremente cerca
de la superficie terrestre pueden tener diferentes
velocidades iniciales, pero experimentar la misma
aceleración gravitatoria.

Despreciando el aire:

\[
\vect{a}_A=-g\jhat,
\]

\[
\vect{a}_B=-g\jhat.
\]

Por tanto:

\[
\vect{a}_{A/B}
=
-g\jhat-(-g\jhat)
=
\vect{0}.
\]

\begin{interpretacion}

Si ambos cuerpos experimentan exactamente la misma
aceleración, esa aceleración no modifica su velocidad
relativa.

El cambio de velocidad que ambos experimentan
es igual y se cancela al restar los vectores.

\end{interpretacion}


\section{Sistemas de referencia móviles}

Hasta ahora hemos comparado directamente dos partículas
respecto de un mismo sistema.

Consideraremos ahora un objeto observado desde dos sistemas
de referencia diferentes.

Sean:

\begin{itemize}
    \item \(S\): sistema asociado al suelo;
    \item \(S'\): sistema que se traslada respecto de \(S\);
    \item \(P\): partícula cuyo movimiento queremos describir.
\end{itemize}

Los orígenes de ambos sistemas pueden estar en posiciones
diferentes.

Denotamos mediante:

\[
\vect{R}(t)
\]

el vector que va desde el origen de \(S\) hasta el origen
de \(S'\), medido en \(S\).

Supondremos que los ejes de ambos sistemas mantienen
orientaciones paralelas.

\begin{center}
\begin{tikzpicture}[scale=0.95]

    \coordinate (O) at (0,0);
    \coordinate (Op) at (3.0,1.1);
    \coordinate (P) at (5.8,3.0);

    \draw[-{Latex},thick]
        (-0.5,0) -- (6.8,0)
        node[right] {\(x\)};

    \draw[-{Latex},thick]
        (0,-0.4) -- (0,4.0)
        node[above] {\(y\)};

    \draw[-{Latex},thick,AzulSecundario]
        (2.5,1.1) -- (6.5,1.1)
        node[right] {\(x'\)};

    \draw[-{Latex},thick,AzulSecundario]
        (3.0,0.7) -- (3.0,4.0)
        node[above] {\(y'\)};

    \draw[-{Latex},very thick,VerdePrincipal]
        (O) -- (Op)
        node[midway,below] {\(\vect{R}\)};

    \draw[-{Latex},very thick,RojoPrincipal]
        (Op) -- (P)
        node[midway,below right] {\(\vect{r}'_P\)};

    \draw[-{Latex},very thick,AzulPrincipal]
        (O) -- (P)
        node[midway,above left] {\(\vect{r}_P\)};

    \fill (O) circle (2pt);
    \fill (Op) circle (2.5pt);
    \fill (P) circle (2.5pt);

    \node[below left] at (O) {\(O\)};
    \node[below right] at (Op) {\(O'\)};
    \node[above right] at (P) {\(P\)};

\end{tikzpicture}
\end{center}


\section{Transformación de posiciones}

La posición de \(P\) respecto de \(S\) se obtiene sumando:

\begin{itemize}
    \item la posición del origen de \(S'\) respecto de \(S\);
    \item la posición de \(P\) respecto de \(S'\).
\end{itemize}

Geométricamente:

\[
\boxed{
\vect{r}_P=\vect{R}+\vect{r}'_P
}.
\]

Esta ecuación expresa que podemos llegar desde el origen
de \(S\) hasta \(P\) siguiendo dos desplazamientos sucesivos.

Si:

\[
\vect{R}=X\ihat+Y\jhat,
\]

y:

\[
\vect{r}'_P=x'\ihat+y'\jhat,
\]

entonces:

\[
\boxed{
x=X+x'
},
\]

\[
\boxed{
y=Y+y'
}.
\]

En tres dimensiones:

\[
\boxed{
z=Z+z'
}.
\]

\begin{interpretacion}

La transformación de posiciones no altera la ubicación
física de la partícula.

Lo que cambia son las coordenadas utilizadas para
describir esa ubicación.

\end{interpretacion}


\section{Transformación de velocidades}

Derivamos la relación:

\[
\vect{r}_P=\vect{R}+\vect{r}'_P.
\]

Como los ejes permanecen paralelos:

\[
\deriv{\vect{r}_P}{t}
=
\deriv{\vect{R}}{t}
+
\deriv{\vect{r}'_P}{t}.
\]

Definimos:

\[
\vect{v}_{P/S}
=
\deriv{\vect{r}_P}{t},
\]

\[
\vect{v}_{S'/S}
=
\deriv{\vect{R}}{t},
\]

\[
\vect{v}_{P/S'}
=
\deriv{\vect{r}'_P}{t}.
\]

Así obtenemos:

\[
\boxed{
\vect{v}_{P/S}
=
\vect{v}_{P/S'}
+
\vect{v}_{S'/S}
}.
\]

Esta relación también se conoce como
\textbf{composición galileana de velocidades}.

\begin{principio}[Composición clásica de velocidades]

Para sistemas de referencia cuyos ejes permanecen paralelos,
la velocidad de una partícula respecto de un sistema puede
obtenerse sumando vectorialmente:

\begin{itemize}
    \item la velocidad de la partícula respecto del sistema móvil;
    \item la velocidad del sistema móvil respecto del sistema
    de referencia considerado.
\end{itemize}

\[
\boxed{
\vect{v}_{P/S}
=
\vect{v}_{P/S'}
+
\vect{v}_{S'/S}
}.
\]

Esta composición corresponde a la cinemática clásica.

\end{principio}


\section{Interpretación de la composición de velocidades}

La ecuación:

\[
\vect{v}_{P/S}
=
\vect{v}_{P/S'}
+
\vect{v}_{S'/S}
\]

puede leerse como:

\[
\text{velocidad respecto del suelo}
\]

igual a:

\[
\text{velocidad respecto del vehículo}
\]

más:

\[
\text{velocidad del vehículo respecto del suelo}.
\]

Por ejemplo, una persona camina dentro de un tren.

La velocidad de la persona respecto del suelo depende
de dos movimientos simultáneos:

\begin{itemize}
    \item el movimiento de la persona dentro del tren;
    \item el movimiento del tren respecto del suelo.
\end{itemize}

\begin{advertencia}

La ecuación es vectorial.

En dos o tres dimensiones no debemos sumar automáticamente
las magnitudes de las velocidades.

Primero debemos sumar sus componentes rectangulares.

\end{advertencia}


\section{Composición de velocidades en una dimensión}

Si todos los movimientos ocurren a lo largo de un mismo eje:

\[
\boxed{
v_{P/S,x}
=
v_{P/S',x}
+
v_{S'/S,x}
}.
\]

Los signos se determinan respecto del eje elegido.

Por ejemplo, tomemos \(+x\) hacia la parte delantera
de un tren.

Si el tren avanza respecto del suelo con:

\[
v_{T/S}=+\SI{15}{\meter\per\second},
\]

y una persona camina hacia adelante con:

\[
v_{P/T}=+\SI{2}{\meter\per\second},
\]

entonces:

\[
v_{P/S}=2+15.
\]

Así:

\[
\boxed{
v_{P/S}=+\SI{17}{\meter\per\second}
}.
\]

Si la persona camina hacia atrás con:

\[
v_{P/T}=-\SI{2}{\meter\per\second},
\]

tenemos:

\[
v_{P/S}=-2+15.
\]

Por tanto:

\[
\boxed{
v_{P/S}=+\SI{13}{\meter\per\second}
}.
\]

\begin{interpretacion}

Caminar hacia atrás dentro del tren no significa
necesariamente desplazarse hacia atrás respecto del suelo.

Aunque la persona retroceda respecto del tren,
puede continuar avanzando respecto del suelo.

\end{interpretacion}


\section{Ejemplo: persona sobre una banda transportadora}

\begin{ejemplo}[Movimiento sobre una banda móvil]

Una banda transportadora horizontal se mueve hacia
la derecha con velocidad:

\[
\vect{v}_{B/S}
=
3.0\ihat\,
\si{\meter\per\second}.
\]

Una persona camina sobre ella hacia la izquierda
con velocidad relativa a la banda:

\[
\vect{v}_{P/B}
=
-1.2\ihat\,
\si{\meter\per\second}.
\]

Determine la velocidad de la persona respecto del suelo.

\end{ejemplo}

Elegimos \(+x\) hacia la derecha.

Aplicamos:

\[
\vect{v}_{P/S}
=
\vect{v}_{P/B}
+
\vect{v}_{B/S}.
\]

Sustituyendo:

\[
\vect{v}_{P/S}
=
(-1.2\ihat+3.0\ihat)
\,\si{\meter\per\second}.
\]

Por tanto:

\[
\boxed{
\vect{v}_{P/S}
=
1.8\ihat\,
\si{\meter\per\second}
}.
\]

\begin{interpretacion}

La persona camina hacia la izquierda respecto de la banda,
pero se mueve hacia la derecha respecto del suelo.

Esto ocurre porque la banda transportadora avanza hacia
la derecha más rápidamente de lo que la persona camina
hacia la izquierda respecto de ella.

\end{interpretacion}


\section{Composición de velocidades en dos dimensiones}

Cuando los movimientos no ocurren sobre una misma recta,
debemos utilizar componentes rectangulares.

La relación vectorial es:

\[
\vect{v}_{P/S}
=
\vect{v}_{P/S'}
+
\vect{v}_{S'/S}.
\]

Si:

\[
\vect{v}_{P/S'}
=
v_{P/S',x}\ihat+v_{P/S',y}\jhat,
\]

y:

\[
\vect{v}_{S'/S}
=
v_{S'/S,x}\ihat+v_{S'/S,y}\jhat,
\]

entonces:

\[
\boxed{
v_{P/S,x}
=
v_{P/S',x}+v_{S'/S,x}
},
\]

\[
\boxed{
v_{P/S,y}
=
v_{P/S',y}+v_{S'/S,y}
}.
\]

Finalmente:

\[
v_{P/S}
=
\sqrt{
v_{P/S,x}^2+v_{P/S,y}^2
}.
\]

La dirección se determina a partir de los signos de
las componentes y del cuadrante correspondiente.


\section{Ejemplo: velocidades perpendiculares}

\begin{ejemplo}[Movimiento sobre una plataforma móvil]

Una plataforma se desplaza hacia \(+x\) con velocidad:

\[
\vect{v}_{B/S}
=
4.0\ihat\,
\si{\meter\per\second}.
\]

Sobre ella, una persona camina hacia \(+y\) con
velocidad relativa:

\[
\vect{v}_{P/B}
=
3.0\jhat\,
\si{\meter\per\second}.
\]

Determine la velocidad de la persona respecto del suelo,
su rapidez y su dirección.

\end{ejemplo}

Aplicamos la composición vectorial:

\[
\vect{v}_{P/S}
=
\vect{v}_{P/B}+\vect{v}_{B/S}.
\]

Sumamos por componentes:

\[
v_{P/S,x}=\SI{4.0}{\meter\per\second},
\]

\[
v_{P/S,y}=\SI{3.0}{\meter\per\second}.
\]

Por tanto:

\[
\boxed{
\vect{v}_{P/S}
=
(4.0\ihat+3.0\jhat)
\,\si{\meter\per\second}
}.
\]

La rapidez es:

\[
v_{P/S}
=
\sqrt{4.0^2+3.0^2}\,
\si{\meter\per\second}.
\]

Así:

\[
\boxed{
v_{P/S}=\SI{5.0}{\meter\per\second}
}.
\]

Para determinar la dirección:

\[
\tan\theta=\frac{3.0}{4.0}.
\]

Por tanto:

\[
\boxed{
\theta\approx36.9\grado
}
\]

por encima de \(+x\).

\begin{interpretacion}

La persona no se mueve exclusivamente en la dirección
en la que camina respecto de la plataforma.

Para el observador del suelo, su movimiento es el resultado
de combinar ambas velocidades.

\end{interpretacion}


\section{Movimiento relativo en tres dimensiones}

La composición de velocidades no cambia conceptualmente
cuando añadimos una tercera dimensión.

Tenemos:

\[
\boxed{
\vect{v}_{P/S}
=
\vect{v}_{P/S'}
+
\vect{v}_{S'/S}
}.
\]

La única diferencia es que ahora necesitamos tres componentes:

\[
\boxed{
v_{P/S,x}
=
v_{P/S',x}+v_{S'/S,x}
},
\]

\[
\boxed{
v_{P/S,y}
=
v_{P/S',y}+v_{S'/S,y}
},
\]

\[
\boxed{
v_{P/S,z}
=
v_{P/S',z}+v_{S'/S,z}
}.
\]

La rapidez es:

\[
\boxed{
v_{P/S}
=
\sqrt{
v_{P/S,x}^2+
v_{P/S,y}^2+
v_{P/S,z}^2
}
}.
\]

\begin{idea}

El movimiento relativo en una, dos o tres dimensiones
se describe mediante el mismo principio vectorial.

La diferencia está en la cantidad de componentes necesarias.

\end{idea}


\section{Transformación de aceleraciones entre sistemas móviles}

Derivamos la composición de velocidades:

\[
\vect{v}_{P/S}
=
\vect{v}_{P/S'}
+
\vect{v}_{S'/S}.
\]

Entonces:

\[
\boxed{
\vect{a}_{P/S}
=
\vect{a}_{P/S'}
+
\vect{a}_{S'/S}
}.
\]

Esta relación es válida para sistemas cuyos ejes
permanecen paralelos y sin rotación relativa.

El término:

\[
\vect{a}_{S'/S}
\]

es la aceleración del origen del sistema móvil
respecto del sistema \(S\).


\section{Sistemas móviles con velocidad constante}

Si el sistema móvil se traslada con velocidad vectorial
constante:

\[
\vect{v}_{S'/S}=\text{constante},
\]

entonces:

\[
\vect{a}_{S'/S}=\vect{0}.
\]

Por tanto:

\[
\boxed{
\vect{a}_{P/S}=\vect{a}_{P/S'}
}.
\]

\begin{principio}[Invariancia galileana de la aceleración]

Dos sistemas cartesianos cuyos ejes permanecen paralelos
y que se mueven uno respecto del otro con velocidad
vectorial constante miden la misma aceleración para
una partícula.

\[
\boxed{
\vect{a}'=\vect{a}
}.
\]

\end{principio}

\begin{interpretacion}

Un cambio entre sistemas de referencia en traslación
uniforme puede modificar la velocidad atribuida
a un objeto.

Sin embargo, no modifica su aceleración clásica.

Esta propiedad será especialmente importante al
estudiar los sistemas de referencia inerciales
en dinámica.

\end{interpretacion}


\section{Sistemas móviles acelerados}

Si el sistema móvil tiene aceleración respecto del suelo:

\[
\vect{a}_{S'/S}\neq\vect{0},
\]

las aceleraciones medidas desde ambos sistemas
no son necesariamente iguales.

De:

\[
\vect{a}_{P/S}
=
\vect{a}_{P/S'}
+
\vect{a}_{S'/S},
\]

obtenemos:

\[
\boxed{
\vect{a}_{P/S'}
=
\vect{a}_{P/S}
-
\vect{a}_{S'/S}
}.
\]

\begin{interpretacion}

Un objeto puede encontrarse en reposo respecto de
un vehículo que acelera, pero poseer aceleración
respecto del suelo.

No existe contradicción.

Las aceleraciones corresponden a sistemas de referencia
diferentes.

\end{interpretacion}


\section{Ejemplo: objeto dentro de un vehículo acelerado}

\begin{ejemplo}

Un automóvil acelera horizontalmente hacia la derecha con:

\[
\vect{a}_{A/S}
=
2.0\ihat\,
\si{\meter\per\second\squared}.
\]

Un objeto permanece fijo respecto del interior
del automóvil.

Determine su aceleración respecto del automóvil
y respecto del suelo.

\end{ejemplo}

Como el objeto permanece fijo respecto del automóvil:

\[
\vect{v}_{P/A}=\vect{0},
\]

y:

\[
\boxed{
\vect{a}_{P/A}=\vect{0}
}.
\]

La composición de aceleraciones es:

\[
\vect{a}_{P/S}
=
\vect{a}_{P/A}
+
\vect{a}_{A/S}.
\]

Sustituyendo:

\[
\vect{a}_{P/S}
=
\vect{0}
+
2.0\ihat\,
\si{\meter\per\second\squared}.
\]

Por tanto:

\[
\boxed{
\vect{a}_{P/S}
=
2.0\ihat\,
\si{\meter\per\second\squared}
}.
\]

\begin{interpretacion}

El objeto está en reposo respecto del automóvil
porque mantiene constante su posición dentro de él.

Sin embargo, acelera respecto del suelo porque
participa del movimiento acelerado del vehículo.

\end{interpretacion}


\section{Movimiento relativo de una embarcación}

Uno de los ejemplos más importantes del movimiento
relativo bidimensional es el de una embarcación
que cruza un río.

En esta situación intervienen tres elementos:

\begin{itemize}
    \item la embarcación;
    \item el agua;
    \item la orilla o el suelo.
\end{itemize}

La velocidad de la embarcación respecto del suelo
es la suma de:

\[
\vect{v}_{B/A},
\]

velocidad de la embarcación respecto del agua, y:

\[
\vect{v}_{A/S},
\]

velocidad del agua respecto del suelo.

Por tanto:

\[
\boxed{
\vect{v}_{B/S}
=
\vect{v}_{B/A}
+
\vect{v}_{A/S}
}.
\]

\begin{advertencia}

La rapidez que desarrolla la embarcación respecto
del agua no es necesariamente igual a su rapidez
respecto de la orilla.

La corriente modifica la velocidad resultante
medida desde el suelo.

\end{advertencia}


\section{Elección de ejes para cruzar un río}

Consideremos un río cuyas orillas son aproximadamente
rectas y paralelas.

Elegimos:

\begin{itemize}
    \item \(+x\): perpendicular a las orillas, hacia la opuesta;
    \item \(+y\): paralelo a la corriente, aguas abajo.
\end{itemize}

La corriente tiene velocidad:

\[
\vect{v}_{A/S}
=
u\jhat,
\]

donde:

\[
u>0.
\]

La velocidad de la embarcación respecto del agua
puede tener componentes:

\[
\vect{v}_{B/A}
=
v_{Bx/A}\ihat+v_{By/A}\jhat.
\]

Entonces:

\[
\boxed{
\vect{v}_{B/S}
=
v_{Bx/A}\ihat
+
(v_{By/A}+u)\jhat
}.
\]

Por componentes:

\[
\boxed{
v_{Bx/S}=v_{Bx/A}
},
\]

\[
\boxed{
v_{By/S}=v_{By/A}+u
}.
\]

Estas dos relaciones permiten analizar diferentes
estrategias de cruce.


\section{Embarcación dirigida perpendicularmente a la orilla}

Supongamos que la embarcación apunta directamente
hacia la orilla opuesta.

Su velocidad respecto del agua es:

\[
\vect{v}_{B/A}=v\ihat.
\]

La corriente lleva:

\[
\vect{v}_{A/S}=u\jhat.
\]

Por tanto:

\[
\boxed{
\vect{v}_{B/S}=v\ihat+u\jhat
}.
\]

La trayectoria observada desde la orilla
no es perpendicular al río.

La embarcación avanza hacia la otra orilla
y simultáneamente deriva aguas abajo.

\begin{center}
\begin{tikzpicture}[scale=0.95]

    \fill[AzulSecundario!10]
        (0,0) rectangle (7.2,4.3);

    \draw[thick]
        (0,0) -- (7.2,0);

    \draw[thick]
        (0,4.3) -- (7.2,4.3);

    \node[above] at (3.6,4.3) {orilla opuesta};

    \node[below] at (3.6,0) {orilla de partida};

    \coordinate (P) at (1.2,0.5);

    \draw[-{Latex},very thick,AzulPrincipal]
        (P) -- ++(0,2.2)
        node[left] {\(\vect{v}_{B/A}\)};

    \draw[-{Latex},very thick,VerdePrincipal]
        (P) -- ++(2.2,0)
        node[below] {\(\vect{v}_{A/S}\)};

    \draw[-{Latex},very thick,RojoPrincipal]
        (P) -- ++(2.2,2.2)
        node[above right] {\(\vect{v}_{B/S}\)};

    \draw[-{Latex},thick,dashed]
        (P) -- ++(4.0,4.0);

    \draw[-{Latex},thick]
        (5.7,2.1) -- (6.7,2.1)
        node[right] {corriente};

\end{tikzpicture}
\end{center}

El dibujo representa la composición vectorial;
la dirección del cruce se ha dispuesto verticalmente
para facilitar su lectura.


\section{Tiempo de cruce}

Sea \(D\) el ancho del río.

Si la componente perpendicular a las orillas
es constante:

\[
v_{Bx/S}=v_{Bx/A},
\]

el tiempo de cruce es:

\[
\boxed{
t=\frac{D}{v_{Bx/A}}
}.
\]

En el caso de una embarcación dirigida
perpendicularmente:

\[
v_{Bx/A}=v.
\]

Entonces:

\[
\boxed{
t=\frac{D}{v}
}.
\]

\begin{interpretacion}

Si la corriente es paralela a las orillas,
no aporta componente de velocidad perpendicular.

Por tanto, para una embarcación que apunta
perpendicularmente, el tiempo de cruce
no depende de la rapidez de la corriente.

La corriente sí modifica el desplazamiento
aguas abajo.

\end{interpretacion}


\section{Deriva aguas abajo}

La velocidad paralela a las orillas es:

\[
v_{By/S}=u.
\]

Si la corriente es uniforme, el desplazamiento
aguas abajo durante el cruce es:

\[
\Delta y=ut.
\]

Sustituyendo:

\[
t=\frac{D}{v},
\]

obtenemos:

\[
\boxed{
\Delta y=\frac{uD}{v}
}.
\]

\begin{condicion}

Esta relación corresponde al caso en el que:

\begin{itemize}
    \item la embarcación apunta perpendicularmente a las orillas;
    \item su velocidad respecto del agua es constante;
    \item la corriente es uniforme y paralela a las orillas;
    \item el río tiene ancho constante en la zona considerada.
\end{itemize}

\end{condicion}


\section{Ejemplo: cruce con deriva}

\begin{ejemplo}

Una embarcación cruza un río de ancho:

\[
D=\SI{120}{\meter}.
\]

Su rapidez respecto del agua es:

\[
v=\SI{4.0}{\meter\per\second},
\]

y apunta perpendicularmente a las orillas.

La corriente tiene rapidez:

\[
u=\SI{3.0}{\meter\per\second}.
\]

Determine:

\begin{enumerate}[label=\alph*)]
    \item El vector velocidad respecto del suelo.
    \item La rapidez respecto del suelo.
    \item El tiempo de cruce.
    \item La deriva aguas abajo.
\end{enumerate}

\end{ejemplo}

Elegimos \(+x\) hacia la orilla opuesta
y \(+y\) aguas abajo.

Las velocidades son:

\[
\vect{v}_{B/A}
=
4.0\ihat\,
\si{\meter\per\second},
\]

\[
\vect{v}_{A/S}
=
3.0\jhat\,
\si{\meter\per\second}.
\]

Entonces:

\[
\boxed{
\vect{v}_{B/S}
=
(4.0\ihat+3.0\jhat)
\,\si{\meter\per\second}
}.
\]

La rapidez respecto del suelo es:

\[
v_{B/S}
=
\sqrt{4.0^2+3.0^2}\,
\si{\meter\per\second}.
\]

Por tanto:

\[
\boxed{
v_{B/S}=\SI{5.0}{\meter\per\second}
}.
\]

El tiempo de cruce depende de la componente
perpendicular:

\[
t=\frac{D}{v_{Bx/S}}.
\]

Entonces:

\[
t=\frac{120}{4.0}.
\]

Así:

\[
\boxed{
t=\SI{30}{\second}
}.
\]

La deriva es:

\[
\Delta y=ut.
\]

Por tanto:

\[
\Delta y=(3.0)(30).
\]

Así:

\[
\boxed{
\Delta y=\SI{90}{\meter}
}.
\]

\begin{interpretacion}

Aunque la embarcación apunta directamente
hacia la orilla opuesta, llega
\(\SI{90}{\meter}\) aguas abajo
del punto situado frente a su posición inicial.

La dirección de la embarcación respecto del agua
no coincide con la dirección de su trayectoria
respecto del suelo.

\end{interpretacion}


\section{Cruzar sin deriva}

Ahora queremos que la embarcación llegue al punto
directamente opuesto al de partida.

Para ello, su desplazamiento paralelo a la corriente
debe ser cero.

Si las velocidades son constantes:

\[
v_{By/S}=0.
\]

Pero:

\[
v_{By/S}=v_{By/A}+u.
\]

Entonces:

\[
v_{By/A}+u=0.
\]

Por tanto:

\[
\boxed{
v_{By/A}=-u
}.
\]

La embarcación debe apuntar parcialmente
aguas arriba para compensar la corriente.

Si su rapidez respecto del agua es \(v\),
sus componentes deben satisfacer:

\[
v^2=v_{Bx/A}^2+v_{By/A}^2.
\]

Como:

\[
v_{By/A}=-u,
\]

tenemos:

\[
v_{Bx/A}^2=v^2-u^2.
\]

Por tanto:

\[
\boxed{
v_{Bx/A}=\sqrt{v^2-u^2}
}.
\]

Elegimos la raíz positiva porque queremos avanzar
hacia la orilla opuesta.

\begin{condicion}

Para cruzar sin deriva con una velocidad de embarcación
constante respecto del agua es necesario:

\[
v>u.
\]

Si:

\[
v=u,
\]

la única manera de cancelar completamente la corriente
es apuntar aguas arriba sin componente de avance
hacia la otra orilla.

Si:

\[
v<u,
\]

no es posible cancelar la corriente y simultáneamente
cruzar bajo este modelo.

\end{condicion}


\section{Ángulo necesario para evitar la deriva}

Sea \(\phi\) el ángulo que forma la dirección de
la embarcación con la perpendicular a las orillas,
medido hacia aguas arriba.

Las componentes son:

\[
v_{Bx/A}=v\cos\phi,
\]

\[
v_{By/A}=-v\sen\phi.
\]

Para evitar la deriva:

\[
-v\sen\phi+u=0.
\]

Por tanto:

\[
\boxed{
\sen\phi=\frac{u}{v}
}.
\]

Así:

\[
\boxed{
\phi=\arcsen\left(\frac{u}{v}\right)
}.
\]

El tiempo de cruce será:

\[
t=\frac{D}{v\cos\phi}.
\]

Utilizando:

\[
v\cos\phi=\sqrt{v^2-u^2},
\]

obtenemos:

\[
\boxed{
t=\frac{D}{\sqrt{v^2-u^2}}
}.
\]


\section{Ejemplo: embarcación sin deriva}

\begin{ejemplo}

Una embarcación debe cruzar un río de:

\[
D=\SI{150}{\meter}
\]

y llegar directamente al punto opuesto.

La corriente tiene rapidez:

\[
u=\SI{3.0}{\meter\per\second}.
\]

La embarcación mantiene una rapidez respecto del agua de:

\[
v=\SI{5.0}{\meter\per\second}.
\]

Determine:

\begin{enumerate}[label=\alph*)]
    \item La dirección que debe adoptar.
    \item Las componentes de su velocidad respecto del agua.
    \item Su velocidad respecto del suelo.
    \item El tiempo de cruce.
\end{enumerate}

\end{ejemplo}

Elegimos \(+x\) perpendicular a las orillas
y \(+y\) aguas abajo.

Para no derivar:

\[
v_{By/S}=0.
\]

Por tanto:

\[
v_{By/A}=-u.
\]

Así:

\[
\boxed{
v_{By/A}=-\SI{3.0}{\meter\per\second}
}.
\]

La componente perpendicular es:

\[
v_{Bx/A}=\sqrt{v^2-u^2}.
\]

Sustituyendo:

\[
v_{Bx/A}=\sqrt{5.0^2-3.0^2}.
\]

Por tanto:

\[
\boxed{
v_{Bx/A}=\SI{4.0}{\meter\per\second}
}.
\]

El vector velocidad respecto del agua es:

\[
\boxed{
\vect{v}_{B/A}
=
(4.0\ihat-3.0\jhat)
\,\si{\meter\per\second}
}.
\]

La dirección respecto de la perpendicular satisface:

\[
\tan\phi=\frac{3.0}{4.0}.
\]

Por tanto:

\[
\boxed{
\phi\approx36.9\grado
}
\]

hacia aguas arriba.

La corriente posee:

\[
\vect{v}_{A/S}
=
3.0\jhat\,
\si{\meter\per\second}.
\]

Sumando:

\[
\vect{v}_{B/S}
=
(4.0\ihat-3.0\jhat)
+
3.0\jhat.
\]

Así:

\[
\boxed{
\vect{v}_{B/S}
=
4.0\ihat\,
\si{\meter\per\second}
}.
\]

La velocidad resultante es perpendicular
a las orillas.

Finalmente:

\[
t=\frac{D}{v_{Bx/S}}.
\]

Entonces:

\[
t=\frac{150}{4.0}.
\]

Por tanto:

\[
\boxed{
t=\SI{37.5}{\second}
}.
\]

\begin{interpretacion}

La embarcación debe orientar su movimiento
aguas arriba, aunque su desplazamiento respecto
del suelo será directamente hacia la orilla opuesta.

Esto muestra que dirección de avance respecto
del agua y trayectoria respecto del suelo
son conceptos distintos.

\end{interpretacion}


\section{Cruce más rápido y cruce sin deriva}

Las dos estrategias estudiadas responden
a objetivos diferentes.

Si la embarcación mantiene rapidez \(v\)
respecto del agua y la corriente es uniforme,
el menor tiempo para alcanzar la orilla opuesta
se obtiene maximizando la componente perpendicular:

\[
v_{Bx/A}.
\]

Esto ocurre cuando la embarcación apunta
perpendicularmente a las orillas:

\[
v_{Bx/A}=v.
\]

Entonces:

\[
\boxed{
t_{\min}=\frac{D}{v}
}.
\]

Pero esa estrategia produce deriva
cuando:

\[
u\neq0.
\]

Para llegar directamente al punto opuesto,
la embarcación debe orientar parte de su
velocidad aguas arriba.

Entonces:

\[
v_{Bx/A}<v.
\]

Por tanto:

\[
t_{\mathrm{sin\ deriva}}>t_{\min}.
\]

\begin{idea}

El trayecto de menor tiempo y el trayecto
que termina directamente frente al punto
de partida no tienen por qué coincidir.

La estrategia depende de qué condición
física queremos satisfacer.

\end{idea}


\section{Movimiento relativo de una aeronave}

El movimiento de una aeronave respecto del suelo
también depende del movimiento del medio.

La aeronave se desplaza respecto del aire,
pero el aire puede moverse respecto del suelo.

Por tanto:

\[
\boxed{
\vect{v}_{P/S}
=
\vect{v}_{P/A}
+
\vect{v}_{A/S}
}.
\]

Aquí:

\[
\vect{v}_{P/A}
\]

representa la velocidad de la aeronave
respecto del aire.

\[
\vect{v}_{A/S}
\]

representa la velocidad del viento respecto del suelo.

\[
\vect{v}_{P/S}
\]

representa la velocidad de la aeronave
respecto del suelo.

\begin{interpretacion}

Una aeronave puede apuntar hacia una dirección
y desplazarse respecto del suelo en otra.

El viento modifica la trayectoria observada
desde la superficie terrestre.

\end{interpretacion}


\section{Viento de cola y viento de frente}

Si todos los movimientos son colineales,
basta trabajar con una componente.

Elegimos \(+x\) en la dirección del vuelo.

Si la aeronave vuela respecto del aire con:

\[
v_{P/A}=+\SI{70}{\meter\per\second},
\]

y existe viento de cola:

\[
v_{A/S}=+\SI{15}{\meter\per\second},
\]

entonces:

\[
v_{P/S}=70+15.
\]

Por tanto:

\[
\boxed{
v_{P/S}=+\SI{85}{\meter\per\second}
}.
\]

Si el viento es de frente:

\[
v_{A/S}=-\SI{15}{\meter\per\second},
\]

entonces:

\[
v_{P/S}=70-15.
\]

Por tanto:

\[
\boxed{
v_{P/S}=+\SI{55}{\meter\per\second}
}.
\]

\begin{interpretacion}

La rapidez de la aeronave respecto del aire
puede mantenerse constante mientras
su rapidez respecto del suelo cambia
debido a la velocidad del viento.

\end{interpretacion}


\section{Viento lateral y corrección de rumbo}

Supongamos que una aeronave debe desplazarse
hacia el norte sin desviarse hacia el este.

Elegimos:

\[
+x:\text{ este},
\]

\[
+y:\text{ norte}.
\]

El viento sopla hacia el este con rapidez \(u\):

\[
\vect{v}_{A/S}=u\ihat.
\]

La aeronave debe apuntar parcialmente
hacia el oeste para compensarlo.

Su velocidad respecto del aire tendrá:

\[
v_{Px/A}<0.
\]

Para que la trayectoria respecto del suelo
no tenga componente horizontal:

\[
v_{Px/S}=0.
\]

Entonces:

\[
v_{Px/A}+u=0.
\]

Por tanto:

\[
\boxed{
v_{Px/A}=-u
}.
\]

Si la rapidez respecto del aire es \(v\),
la componente hacia el norte será:

\[
\boxed{
v_{Py/A}=\sqrt{v^2-u^2}
}.
\]

De esta manera, el viento lateral
se compensa mediante una orientación
apropiada de la aeronave.

\begin{condicion}

Esta construcción supone velocidades constantes,
ejes paralelos y un viento uniforme durante
el intervalo considerado.

También requiere:

\[
v>u
\]

para conservar una componente de avance
hacia el norte después de compensar el viento.

\end{condicion}


\section{Ejemplo: aeronave con viento lateral}

\begin{ejemplo}

Una aeronave debe desplazarse hacia el norte.

Su rapidez respecto del aire es:

\[
v=\SI{80.0}{\meter\per\second}.
\]

El viento sopla hacia el este con rapidez:

\[
u=\SI{20.0}{\meter\per\second}.
\]

Determine el ángulo de corrección que debe adoptar
para seguir una trayectoria directamente hacia el norte
y su rapidez respecto del suelo.

\end{ejemplo}

Elegimos \(+x\) hacia el este y \(+y\)
hacia el norte.

El viento tiene velocidad:

\[
\vect{v}_{A/S}
=
20.0\ihat\,
\si{\meter\per\second}.
\]

Para anular la componente este de la
velocidad respecto del suelo:

\[
v_{Px/A}=-\SI{20.0}{\meter\per\second}.
\]

Como la magnitud de la velocidad
respecto del aire es \(\SI{80.0}{\meter\per\second}\):

\[
v_{Py/A}
=
\sqrt{80.0^2-20.0^2}.
\]

Entonces:

\[
v_{Py/A}
=
\sqrt{6000}\,
\si{\meter\per\second}.
\]

Así:

\[
\boxed{
v_{Py/A}\approx
\SI{77.5}{\meter\per\second}
}.
\]

El ángulo de corrección \(\phi\),
medido hacia el oeste desde el norte,
satisface:

\[
\sen\phi=\frac{20.0}{80.0}.
\]

Por tanto:

\[
\boxed{
\phi\approx14.5\grado
}.
\]

La velocidad de la aeronave respecto del suelo es:

\[
\vect{v}_{P/S}
=
\vect{v}_{P/A}+\vect{v}_{A/S}.
\]

Entonces:

\[
\vect{v}_{P/S}
=
(-20.0\ihat+77.5\jhat)
+
20.0\ihat.
\]

Por tanto:

\[
\boxed{
\vect{v}_{P/S}
=
77.5\jhat\,
\si{\meter\per\second}
}.
\]

\begin{interpretacion}

La aeronave apunta hacia el noroeste,
pero se desplaza respecto del suelo
directamente hacia el norte.

La corrección de rumbo cancela
la componente lateral del viento.

\end{interpretacion}


\section{Movimiento relativo entre dos objetos en el plano}

El movimiento relativo también permite estudiar
la separación entre dos objetos que se desplazan
en diferentes direcciones.

Supongamos que:

\[
\vect{r}_A(t)
\]

y:

\[
\vect{r}_B(t)
\]

son sus posiciones respecto del suelo.

La posición relativa es:

\[
\boxed{
\vect{r}_{A/B}(t)
=
\vect{r}_A(t)-\vect{r}_B(t)
}.
\]

La velocidad relativa es:

\[
\boxed{
\vect{v}_{A/B}(t)
=
\vect{v}_A(t)-\vect{v}_B(t)
}.
\]

Si ambas velocidades son constantes:

\[
\vect{r}_A(t)
=
\vect{r}_{A0}+\vect{v}_At,
\]

\[
\vect{r}_B(t)
=
\vect{r}_{B0}+\vect{v}_Bt.
\]

Restando:

\[
\vect{r}_{A/B}(t)
=
(\vect{r}_{A0}-\vect{r}_{B0})
+
(\vect{v}_A-\vect{v}_B)t.
\]

Por tanto:

\[
\boxed{
\vect{r}_{A/B}(t)
=
\vect{r}_{A/B,0}
+
\vect{v}_{A/B}t
}.
\]

Esta es una ecuación especialmente útil
para estudiar encuentros y aproximaciones.


\section{Condición de encuentro}

Dos partículas se encuentran cuando ocupan
la misma posición en el mismo instante.

Por tanto:

\[
\vect{r}_A(t)=\vect{r}_B(t).
\]

Equivalentemente:

\[
\boxed{
\vect{r}_{A/B}(t)=\vect{0}
}.
\]

En dos dimensiones esto requiere:

\[
\boxed{
x_{A/B}(t)=0
},
\]

\[
\boxed{
y_{A/B}(t)=0
}.
\]

Ambas condiciones deben cumplirse
para el mismo valor de \(t\).

\begin{advertencia}

Que dos partículas tengan una coordenada igual
no significa necesariamente que se encuentren.

En dos dimensiones deben coincidir
simultáneamente sus coordenadas \(x\) e \(y\).

En tres dimensiones también debe coincidir \(z\).

\end{advertencia}


\section{Ejemplo: encuentro en dos dimensiones}

\begin{ejemplo}

Dos partículas se mueven en un plano
con velocidades constantes.

En \(t=0\):

\[
\vect{r}_{A0}
=
(0\ihat+20\jhat)\,\si{\meter},
\]

\[
\vect{r}_{B0}
=
(30\ihat+0\jhat)\,\si{\meter}.
\]

Sus velocidades son:

\[
\vect{v}_A
=
(5\ihat+2\jhat)\,
\si{\meter\per\second},
\]

\[
\vect{v}_B
=
(2\ihat+4\jhat)\,
\si{\meter\per\second}.
\]

Determine si se encuentran
durante su movimiento.

\end{ejemplo}

La posición relativa inicial es:

\[
\vect{r}_{A/B,0}
=
\vect{r}_{A0}-\vect{r}_{B0}.
\]

Por tanto:

\[
\boxed{
\vect{r}_{A/B,0}
=
(-30\ihat+20\jhat)\,\si{\meter}
}.
\]

La velocidad relativa es:

\[
\vect{v}_{A/B}
=
\vect{v}_A-\vect{v}_B.
\]

Entonces:

\[
\boxed{
\vect{v}_{A/B}
=
(3\ihat-2\jhat)\,
\si{\meter\per\second}
}.
\]

La posición relativa es:

\[
\vect{r}_{A/B}(t)
=
(-30+3t)\ihat+(20-2t)\jhat.
\]

Para el encuentro:

\[
-30+3t=0,
\]

\[
20-2t=0.
\]

De la primera:

\[
t=\SI{10}{\second}.
\]

De la segunda:

\[
t=\SI{10}{\second}.
\]

Como ambas condiciones se satisfacen
en el mismo instante:

\[
\boxed{
t=\SI{10}{\second}
}.
\]

La posición de encuentro se obtiene
con la ecuación de \(A\):

\[
\vect{r}_A(t)
=
\vect{r}_{A0}+\vect{v}_At.
\]

Entonces:

\[
\vect{r}_A(10)
=
(0\ihat+20\jhat)
+
(5\ihat+2\jhat)(10).
\]

Por tanto:

\[
\boxed{
\vect{r}_{\mathrm{encuentro}}
=
(50\ihat+40\jhat)\,\si{\meter}
}.
\]

\begin{interpretacion}

Las partículas tienen velocidades distintas
y parten de posiciones diferentes.

Sin embargo, sus movimientos permiten
que ambas ocupen simultáneamente
el punto \((50,40)\,\si{\meter}\).

El método relativo convierte el problema
de dos movimientos en el movimiento
de una posición de separación.

\end{interpretacion}


\section{Distancia entre dos partículas móviles}

La posición relativa permite conocer
la separación entre dos partículas:

\[
\vect{r}_{A/B}
=
\vect{r}_A-\vect{r}_B.
\]

La distancia entre ellas es:

\[
\boxed{
d(t)=\norm{\vect{r}_{A/B}(t)}
}.
\]

En dos dimensiones:

\[
\boxed{
d(t)
=
\sqrt{
x_{A/B}^2(t)+y_{A/B}^2(t)
}
}.
\]

Si la posición relativa depende del tiempo,
la distancia también puede cambiar.

Pero la dirección de la velocidad relativa
no tiene por qué coincidir con la dirección
de la separación.


\section{Rapidez de acercamiento en varias dimensiones}

Definamos:

\[
\vect{r}
=
\vect{r}_{A/B}.
\]

La distancia es:

\[
d=\norm{\vect{r}}.
\]

Entonces:

\[
d^2=\vect{r}\pdot\vect{r}.
\]

Derivando respecto del tiempo:

\[
2d\deriv{d}{t}
=
2\vect{r}\pdot\deriv{\vect{r}}{t}.
\]

Como:

\[
\deriv{\vect{r}}{t}
=
\vect{v}_{A/B},
\]

obtenemos:

\[
\boxed{
\deriv{d}{t}
=
\frac{
\vect{r}_{A/B}\pdot\vect{v}_{A/B}
}{
\norm{\vect{r}_{A/B}}
}
}.
\]

Esta relación es válida cuando:

\[
d\neq0.
\]

\begin{interpretacion}

La tasa de cambio de la distancia depende
de la componente de velocidad relativa
en la dirección de separación.

Si:

\[
\deriv{d}{t}<0,
\]

las partículas se están acercando.

Si:

\[
\deriv{d}{t}>0,
\]

se están alejando.

Si:

\[
\deriv{d}{t}=0,
\]

la distancia no cambia instantáneamente,
aunque ambas partículas puedan continuar
en movimiento.

\end{interpretacion}


\section{Máxima aproximación entre dos partículas}

Consideremos dos partículas que tienen
velocidades constantes.

Definimos:

\[
\vect{r}_0=\vect{r}_{A/B,0},
\]

\[
\vect{v}=\vect{v}_{A/B}.
\]

La separación relativa es:

\[
\vect{r}(t)=\vect{r}_0+\vect{v}t.
\]

La distancia al cuadrado es:

\[
d^2(t)
=
(\vect{r}_0+\vect{v}t)
\pdot
(\vect{r}_0+\vect{v}t).
\]

Desarrollando:

\[
d^2(t)
=
r_0^2
+
2(\vect{r}_0\pdot\vect{v})t
+
v^2t^2.
\]

Para encontrar un mínimo de distancia,
podemos derivar \(d^2\):

\[
\deriv{(d^2)}{t}
=
2\vect{r}_0\pdot\vect{v}
+
2v^2t.
\]

En el instante de mínima separación:

\[
\deriv{(d^2)}{t}=0.
\]

Por tanto:

\[
\vect{r}_0\pdot\vect{v}
+
v^2t_{\min}=0.
\]

Despejando:

\[
\boxed{
t_{\min}
=
-\frac{
\vect{r}_0\pdot\vect{v}
}{
v^2
}
}.
\]

\begin{condicion}

Esta expresión supone velocidad relativa constante
y distinta de cero.

Si sólo estudiamos el intervalo posterior
a \(t=0\), un resultado:

\[
t_{\min}<0
\]

indica que la máxima aproximación de
las trayectorias rectilíneas prolongadas
ocurrió antes del instante inicial.

En ese caso, la mínima distancia del intervalo
\(t\geq0\) ocurre en \(t=0\).

\end{condicion}

Después de determinar el instante físicamente
válido, calculamos:

\[
\boxed{
d_{\min}
=
\norm{\vect{r}_0+\vect{v}t_{\min}}
}.
\]

\begin{interpretacion}

Dos objetos pueden pasar muy cerca sin encontrarse.

La condición de mínima distancia no exige:

\[
\vect{r}_{A/B}=\vect{0}.
\]

Sólo exige que la distancia deje de disminuir
instantáneamente y comience a aumentar,
o que alcance el extremo del intervalo estudiado.

\end{interpretacion}


\section{Ejemplo: distancia mínima entre vehículos}

\begin{ejemplo}

Dos vehículos se desplazan con velocidades constantes
sobre un plano.

En \(t=0\), sus posiciones relativas satisfacen:

\[
\vect{r}_{A/B,0}
=
(100\ihat+60\jhat)\,\si{\meter}.
\]

Su velocidad relativa es:

\[
\vect{v}_{A/B}
=
(-10\ihat-2\jhat)\,
\si{\meter\per\second}.
\]

Determine el instante de máxima aproximación
y la distancia mínima entre ambos.

\end{ejemplo}

Definimos:

\[
\vect{r}_0
=
(100\ihat+60\jhat)\,\si{\meter},
\]

\[
\vect{v}
=
(-10\ihat-2\jhat)\,
\si{\meter\per\second}.
\]

La velocidad relativa es constante.

Calculamos el producto escalar:

\[
\vect{r}_0\pdot\vect{v}
=
100(-10)+60(-2).
\]

Por tanto:

\[
\vect{r}_0\pdot\vect{v}
=
-1120\,
\si{\meter\squared\per\second}.
\]

La magnitud al cuadrado de la velocidad relativa es:

\[
v^2=(-10)^2+(-2)^2.
\]

Así:

\[
v^2
=
104\,\si{\meter\squared\per\second\squared}.
\]

El instante de mínima distancia es:

\[
t_{\min}
=
-\frac{\vect{r}_0\pdot\vect{v}}{v^2}.
\]

Entonces:

\[
t_{\min}=\frac{1120}{104}.
\]

Por tanto:

\[
\boxed{
t_{\min}\approx\SI{10.77}{\second}
}.
\]

La posición relativa en ese instante es:

\[
\vect{r}(t_{\min})
=
(100-10t_{\min})\ihat
+
(60-2t_{\min})\jhat.
\]

Sustituyendo:

\[
x_{A/B}\approx-\SI{7.69}{\meter},
\]

\[
y_{A/B}\approx\SI{38.46}{\meter}.
\]

La distancia mínima es:

\[
d_{\min}
=
\sqrt{
(-7.69)^2+(38.46)^2
}\,\si{\meter}.
\]

Así:

\[
\boxed{
d_{\min}\approx\SI{39.2}{\meter}
}.
\]

\begin{interpretacion}

Los vehículos se acercan inicialmente,
pero no llegan a encontrarse.

Su separación mínima es aproximadamente
\(\SI{39.2}{\meter}\).

Después de ese instante, la distancia
entre ellos aumenta.

\end{interpretacion}


\section{Movimiento relativo y trayectorias}

Una misma partícula puede describir trayectorias
diferentes desde sistemas de referencia distintos.

Consideremos una partícula cuya posición en un sistema
móvil es:

\[
\vect{r}'(t).
\]

Si el origen de ese sistema móvil se desplaza según:

\[
\vect{R}(t),
\]

entonces:

\[
\vect{r}(t)=\vect{R}(t)+\vect{r}'(t).
\]

La trayectoria respecto del suelo se obtiene
combinando ambas funciones del tiempo.

Por ejemplo, una pelota lanzada verticalmente
dentro de un vehículo que se mueve horizontalmente
puede describir:

\begin{itemize}
    \item una trayectoria vertical respecto del vehículo;
    \item una trayectoria parabólica respecto del suelo.
\end{itemize}

Ambas descripciones pueden ser correctas
para el mismo movimiento físico.


\section{Ejemplo: pelota lanzada dentro de un tren}

\begin{ejemplo}

Un tren se mueve horizontalmente hacia \(+x\)
con velocidad constante:

\[
\vect{v}_{T/S}
=
12\ihat\,
\si{\meter\per\second}.
\]

Un pasajero lanza una pelota verticalmente hacia arriba
con velocidad inicial respecto del tren:

\[
\vect{v}_{P/T,0}
=
8\jhat\,
\si{\meter\per\second}.
\]

Despreciando la resistencia del aire, determine:

\begin{enumerate}[label=\alph*)]
    \item La velocidad inicial de la pelota
    respecto del suelo.
    \item Sus ecuaciones de posición respecto del suelo,
    tomando el lanzamiento en el origen.
    \item La forma de la trayectoria vista desde
    el tren y desde el suelo.
\end{enumerate}

\end{ejemplo}

\subsection*{Velocidad inicial}

Aplicamos:

\[
\vect{v}_{P/S,0}
=
\vect{v}_{P/T,0}
+
\vect{v}_{T/S}.
\]

Por tanto:

\[
\vect{v}_{P/S,0}
=
8\jhat+12\ihat.
\]

Así:

\[
\boxed{
\vect{v}_{P/S,0}
=
(12\ihat+8\jhat)\,
\si{\meter\per\second}
}.
\]

\subsection*{Movimiento respecto del suelo}

Tomamos el instante de lanzamiento como \(t=0\).

Las componentes iniciales son:

\[
v_{0x}=\SI{12}{\meter\per\second},
\]

\[
v_{0y}=\SI{8}{\meter\per\second}.
\]

La aceleración es:

\[
\vect{a}_{P/S}=-g\jhat.
\]

Horizontalmente:

\[
\boxed{
x(t)=12t
}.
\]

Verticalmente:

\[
\boxed{
y(t)=8t-\frac12gt^2
}.
\]

Las coordenadas se expresan en metros
cuando \(t\) se expresa en segundos.

Eliminando el tiempo:

\[
t=\frac{x}{12}.
\]

Entonces:

\[
y
=
8\left(\frac{x}{12}\right)
-
\frac12g\left(\frac{x}{12}\right)^2.
\]

Por tanto, la trayectoria respecto del suelo
es una parábola.

\subsection*{Movimiento respecto del tren}

Como el tren se mueve horizontalmente
con velocidad constante, su posición relativa
al suelo es:

\[
X_T(t)=12t.
\]

La coordenada horizontal relativa de la pelota es:

\[
x'(t)=x(t)-X_T(t).
\]

Entonces:

\[
x'(t)=12t-12t=0.
\]

Por tanto:

\[
\boxed{
x'(t)=0
}.
\]

La coordenada vertical es:

\[
\boxed{
y'(t)=8t-\frac12gt^2
}.
\]

Desde el tren, la pelota describe
una trayectoria vertical.

\begin{interpretacion}

Desde el tren, la pelota sube y baja
sobre una misma línea vertical.

Desde el suelo, la pelota avanza
horizontalmente mientras sube y baja.

No se trata de dos movimientos físicos distintos.

Son dos descripciones del mismo movimiento
desde sistemas de referencia diferentes.

\end{interpretacion}


\section{Transformaciones galileanas}

Cuando el sistema \(S'\) se mueve respecto de \(S\)
con velocidad vectorial constante:

\[
\vect{u},
\]

podemos escribir:

\[
\vect{R}(t)=\vect{R}_0+\vect{u}t.
\]

La relación entre posiciones es:

\[
\vect{r}=\vect{R}+\vect{r}'.
\]

Entonces:

\[
\vect{r}'
=
\vect{r}-\vect{R}_0-\vect{u}t.
\]

Si los orígenes coinciden cuando \(t=0\):

\[
\vect{R}_0=\vect{0}.
\]

Así:

\[
\boxed{
\vect{r}'=\vect{r}-\vect{u}t
}.
\]

Estas relaciones forman parte de las
transformaciones galileanas.


\section{Transformaciones por componentes}

Si el sistema móvil se desplaza con:

\[
\vect{u}
=
u_x\ihat+u_y\jhat+u_z\khat,
\]

y los orígenes coinciden en \(t=0\),
entonces:

\[
\boxed{
x'=x-u_xt
},
\]

\[
\boxed{
y'=y-u_yt
},
\]

\[
\boxed{
z'=z-u_zt
}.
\]

Derivando:

\[
\boxed{
v_x'=v_x-u_x
},
\]

\[
\boxed{
v_y'=v_y-u_y
},
\]

\[
\boxed{
v_z'=v_z-u_z
}.
\]

Como las componentes de \(\vect{u}\)
son constantes:

\[
\boxed{
a_x'=a_x
},
\]

\[
\boxed{
a_y'=a_y
},
\]

\[
\boxed{
a_z'=a_z
}.
\]

En la descripción clásica también se supone:

\[
\boxed{
t'=t
}.
\]

Es decir, el tiempo transcurrido es el mismo
para ambos sistemas.


\section{Condiciones de validez de las transformaciones galileanas}

Las transformaciones desarrolladas corresponden
a la mecánica clásica.

Su aplicación supone:

\begin{itemize}
    \item velocidades mucho menores que la velocidad
    de la luz;
    \item ejes cartesianos paralelos;
    \item ausencia de rotación relativa;
    \item traslación uniforme entre los sistemas
    cuando se utiliza la forma galileana
    con \(\vect{u}\) constante.
\end{itemize}

\begin{advertencia}

La expresión:

\[
\vect{v}_{P/S}
=
\vect{v}_{P/S'}+\vect{u}
\]

no es una regla universal para todas las velocidades
y todos los sistemas de referencia.

Cuando las velocidades son comparables
con la velocidad de la luz, debemos utilizar
las transformaciones de la relatividad especial.

Asimismo, cuando los sistemas de referencia giran
uno respecto del otro, aparecen términos adicionales
en las relaciones de velocidad y aceleración.

\end{advertencia}


\section{Sistemas de referencia en rotación}

Hasta ahora hemos supuesto que los ejes del sistema
móvil permanecen paralelos a los ejes del sistema fijo.

Un sistema que gira constituye una situación diferente.

Por ejemplo:

\begin{itemize}
    \item una persona sobre una plataforma giratoria;
    \item un observador dentro de un carrusel;
    \item un sistema de referencia unido a una rueda.
\end{itemize}

En estos casos, los vectores unitarios del sistema
móvil cambian de dirección respecto del sistema fijo.

Por ello, al derivar una posición expresada
en el sistema giratorio, también deben considerarse
las derivadas de esos vectores unitarios.

En general, no basta con escribir:

\[
\vect{v}_{P/S}
=
\vect{v}_{P/S'}
+
\vect{v}_{S'/S}
\]

si los ejes de \(S'\) están girando.

\begin{observacion}

La descripción completa de sistemas de referencia
en rotación requiere incorporar la velocidad angular
del sistema móvil.

En dinámica, además, aparecen las llamadas
fuerzas aparentes asociadas a sistemas
de referencia no inerciales.

Estos efectos deben distinguirse de la simple
composición de velocidades por traslación.

\end{observacion}


\section{Movimiento relativo y sistemas inerciales}

Un sistema de referencia inercial es aquel
en el que se cumplen las leyes de Newton
en su formulación habitual.

En la mecánica clásica, si \(S\) es inercial
y \(S'\) se traslada respecto de \(S\)
con velocidad vectorial constante, sin rotación,
entonces \(S'\) también es inercial.

En ambos sistemas:

\[
\vect{a}'=\vect{a}.
\]

Esta propiedad permite que las leyes
de la dinámica tengan la misma forma
en todos los sistemas inerciales clásicos.

\begin{interpretacion}

La posición y la velocidad dependen
del sistema de referencia.

Sin embargo, la aceleración es la misma
para observadores inerciales relacionados
por transformaciones galileanas.

Esto conecta la cinemática relativa
con el estudio posterior de las leyes
del movimiento.

\end{interpretacion}


\section{Aplicaciones del movimiento relativo}

\begin{aplicacion}[Navegación]

La navegación de embarcaciones requiere distinguir
entre la velocidad respecto del agua
y la velocidad respecto del suelo.

La corriente puede modificar el rumbo efectivo,
el tiempo de llegada y la posición final.

\end{aplicacion}

\begin{aplicacion}[Aviación]

En aviación es necesario distinguir
entre la velocidad respecto del aire
y la velocidad respecto del suelo.

El viento puede alterar la trayectoria
y el tiempo necesario para recorrer
una determinada distancia.

\end{aplicacion}

\begin{aplicacion}[Tránsito vehicular]

La velocidad relativa permite analizar
el acercamiento, adelantamiento y separación
entre vehículos.

También facilita el estudio de encuentros
en una o varias dimensiones.

\end{aplicacion}

\begin{aplicacion}[Seguimiento de objetos]

En sistemas de seguimiento, las posiciones
de diferentes objetos pueden describirse
respecto de un observador o vehículo móvil.

Las diferencias de posición y velocidad
permiten analizar cómo cambia la separación
entre ellos.

\end{aplicacion}


\section{Estrategia general para resolver problemas}

\begin{estrategia}[Movimiento relativo]

Para resolver un problema de movimiento relativo:

\begin{enumerate}

    \item Identifique todos los objetos y sistemas
    de referencia involucrados.

    \item Establezca qué movimiento se desea
    describir y respecto de cuál referencia.

    \item Defina un sistema cartesiano con origen,
    ejes y sentidos positivos.

    \item Dibuje un esquema cuando existan
    varias direcciones de movimiento.

    \item Escriba explícitamente las velocidades
    utilizando subíndices de referencia.

    \item Construya la relación vectorial
    correspondiente:
    \[
    \vect{v}_{A/C}
    =
    \vect{v}_{A/B}
    +
    \vect{v}_{B/C}.
    \]

    \item Descomponga los vectores en componentes
    rectangulares.

    \item Resuelva primero las componentes desconocidas.

    \item Reconstruya el vector resultante.

    \item Determine magnitud y dirección,
    verificando los signos y el cuadrante.

    \item Si se trata de un encuentro, establezca
    la igualdad de posiciones en el mismo instante.

    \item Si se estudia una mínima separación,
    analice la evolución de la posición relativa.

    \item Verifique las condiciones de validez
    del modelo elegido.

    \item Interprete físicamente el resultado.

\end{enumerate}

\end{estrategia}


\section{Errores frecuentes}

\begin{errorfrecuente}[Omitir la referencia]

No es suficiente escribir:

\[
v_A=\SI{20}{\meter\per\second}
\]

si existen varios sistemas de referencia.

Debe indicarse respecto de qué sistema
se mide la velocidad, por ejemplo:

\[
\vect{v}_{A/S}.
\]

\end{errorfrecuente}


\begin{errorfrecuente}[Confundir rapidez con velocidad]

La velocidad relativa es vectorial:

\[
\vect{v}_{A/B}
=
\vect{v}_A-\vect{v}_B.
\]

No debe obtenerse restando magnitudes
sin considerar direcciones y signos.

\end{errorfrecuente}


\begin{errorfrecuente}[Invertir los subíndices]

Las velocidades:

\[
\vect{v}_{A/B}
\]

y:

\[
\vect{v}_{B/A}
\]

tienen sentidos opuestos.

La relación correcta es:

\[
\vect{v}_{A/B}
=
-\vect{v}_{B/A}.
\]

\end{errorfrecuente}


\begin{errorfrecuente}[Confundir dirección de orientación y trayectoria]

Una embarcación o aeronave puede apuntar
hacia una dirección diferente de su
velocidad respecto del suelo.

La trayectoria depende de la velocidad
resultante, no únicamente de la dirección
del vehículo respecto del medio.

\end{errorfrecuente}


\begin{errorfrecuente}[Confundir rapidez relativa y acercamiento]

En dos dimensiones, la magnitud:

\[
\norm{\vect{v}_{A/B}}
\]

no representa necesariamente la rapidez
con la que disminuye la distancia entre
las partículas.

Debe considerarse la dirección de la
velocidad relativa respecto de la
separación.

\end{errorfrecuente}


\begin{errorfrecuente}[Aplicar composición por traslación a sistemas giratorios]

Las relaciones desarrolladas para ejes
paralelos no deben aplicarse directamente
a un sistema de referencia que gira.

En esos casos aparecen contribuciones
adicionales asociadas a la rotación.

\end{errorfrecuente}


\begin{errorfrecuente}[Utilizar tiempos diferentes para un encuentro]

Para que dos partículas se encuentren,
sus posiciones deben coincidir
en el mismo instante.

No basta con que ambas pasen por
el mismo lugar en momentos diferentes.

\end{errorfrecuente}


\section{Secuencia sugerida para impartir la clase}

\begin{notapedagogica}

Una secuencia apropiada para desarrollar
este capítulo es:

\begin{enumerate}

    \item Presentar una persona que camina
    dentro de un tren.

    \item Preguntar si su velocidad es la misma
    para un pasajero y para alguien en la estación.

    \item Recordar el concepto de sistema
    de referencia.

    \item Introducir la notación:
    \[
    \vect{r}_{A/B},
    \qquad
    \vect{v}_{A/B}.
    \]

    \item Construir geométricamente
    la posición relativa entre dos partículas.

    \item Deducir:
    \[
    \vect{v}_{A/B}
    =
    \vect{v}_A-\vect{v}_B.
    \]

    \item Resolver ejemplos rectilíneos
    con igual sentido y sentidos contrarios.

    \item Introducir la composición
    de velocidades entre tres referencias.

    \item Explicar el modelo de sistemas
    con ejes paralelos.

    \item Resolver el ejemplo de banda
    transportadora.

    \item Extender la composición
    a componentes rectangulares.

    \item Resolver un ejemplo bidimensional
    con velocidades perpendiculares.

    \item Analizar una embarcación
    que cruza un río con deriva.

    \item Plantear cómo evitar la deriva
    mediante una orientación aguas arriba.

    \item Comparar cruce más rápido
    y cruce sin deriva.

    \item Aplicar el mismo razonamiento
    al movimiento de una aeronave con viento.

    \item Introducir velocidad y aceleración
    relativas entre dos partículas.

    \item Resolver un problema de encuentro
    en dos dimensiones.

    \item Analizar el concepto
    de distancia mínima.

    \item Presentar una pelota lanzada
    verticalmente dentro de un tren.

    \item Comparar sus trayectorias
    desde ambos sistemas de referencia.

    \item Deducir las transformaciones
    galileanas.

    \item Explicar las condiciones
    de validez de la composición clásica.

    \item Finalizar con preguntas conceptuales
    que distingan movimiento relativo,
    composición vectorial y sistemas
    de referencia.

\end{enumerate}

\end{notapedagogica}


\section{Preguntas de comprensión}

\begin{enumerate}

    \item ¿Por qué la velocidad de un objeto
    depende del sistema de referencia?

    \item ¿Qué significa la expresión
    \(\vect{v}_{A/B}\)?

    \item ¿Qué diferencia existe entre
    \(\vect{v}_{A/B}\) y \(\vect{v}_{B/A}\)?

    \item ¿Puede un objeto encontrarse en reposo
    respecto de un observador y moverse
    respecto de otro?

    \item ¿Qué representa físicamente
    la posición relativa?

    \item ¿Cuál es la diferencia entre
    posición relativa y distancia entre partículas?

    \item ¿Por qué la velocidad relativa
    debe tratarse como un vector?

    \item ¿Cuándo una resta de velocidades
    produce una magnitud igual a la suma
    de las rapideces?

    \item ¿Es suficiente conocer las rapideces
    de dos vehículos para determinar
    su velocidad relativa?

    \item ¿Qué diferencia existe entre
    velocidad relativa y rapidez de acercamiento?

    \item ¿Por qué una embarcación puede
    derivar aunque apunte perpendicularmente
    a la orilla?

    \item ¿De qué depende el tiempo de cruce
    de un río con corriente uniforme?

    \item ¿Cómo puede una embarcación evitar
    la deriva aguas abajo?

    \item ¿Qué condición debe cumplirse
    para cruzar sin deriva manteniendo
    constante la velocidad respecto del agua?

    \item ¿Por qué el cruce más rápido
    no coincide necesariamente con
    el cruce sin deriva?

    \item ¿Qué diferencia existe entre
    la velocidad de una aeronave respecto
    del aire y respecto del suelo?

    \item ¿Qué condición debe cumplirse
    para que dos partículas se encuentren?

    \item ¿Pueden dos partículas pasar
    por el mismo punto y no encontrarse?

    \item ¿Qué significa que la distancia
    entre dos partículas alcance un mínimo?

    \item ¿Puede existir velocidad relativa
    no nula mientras la distancia
    permanece instantáneamente constante?

    \item ¿Por qué dos observadores pueden
    describir trayectorias diferentes
    para un mismo objeto?

    \item ¿Qué condiciones permiten utilizar
    las transformaciones galileanas?

    \item ¿Por qué dos sistemas en traslación
    uniforme miden la misma aceleración?

    \item ¿Qué cambia cuando uno de
    los sistemas de referencia acelera?

    \item ¿Por qué un sistema en rotación
    requiere un tratamiento diferente?

    \item ¿Por qué la composición clásica
    de velocidades no puede aplicarse
    universalmente a velocidades
    comparables con la velocidad de la luz?

\end{enumerate}


\begin{resumen}

El movimiento relativo permite describir
cómo cambia la posición de un objeto
respecto de diferentes sistemas de referencia.

Las posiciones, velocidades y aceleraciones
deben expresarse especificando claramente
la referencia utilizada.

La posición relativa de \(A\)
respecto de \(B\) es:

\[
\boxed{
\vect{r}_{A/B}
=
\vect{r}_A-\vect{r}_B
}.
\]

La velocidad relativa es:

\[
\boxed{
\vect{v}_{A/B}
=
\vect{v}_A-\vect{v}_B
}.
\]

La aceleración relativa es:

\[
\boxed{
\vect{a}_{A/B}
=
\vect{a}_A-\vect{a}_B
}.
\]

Para sistemas cuyos ejes permanecen paralelos:

\[
\boxed{
\vect{r}_{P/S}
=
\vect{R}
+
\vect{r}_{P/S'}
}.
\]

La composición de velocidades es:

\[
\boxed{
\vect{v}_{P/S}
=
\vect{v}_{P/S'}
+
\vect{v}_{S'/S}
}.
\]

Esta relación debe resolverse
mediante componentes rectangulares
cuando intervienen varias dimensiones.

La composición de aceleraciones es:

\[
\boxed{
\vect{a}_{P/S}
=
\vect{a}_{P/S'}
+
\vect{a}_{S'/S}
}.
\]

Si el sistema móvil se traslada
con velocidad vectorial constante:

\[
\boxed{
\vect{a}_{P/S}
=
\vect{a}_{P/S'}
}.
\]

En mecánica clásica, las transformaciones
galileanas relacionan sistemas inerciales
en traslación uniforme.

Si sus orígenes coinciden inicialmente:

\[
\boxed{
\vect{r}'=\vect{r}-\vect{u}t
},
\]

\[
\boxed{
\vect{v}'=\vect{v}-\vect{u}
},
\]

\[
\boxed{
\vect{a}'=\vect{a}
}.
\]

La aplicación del movimiento relativo
permite estudiar:

\begin{itemize}
    \item personas que caminan en vehículos;
    \item automóviles que se persiguen;
    \item embarcaciones en ríos;
    \item aeronaves bajo la influencia del viento;
    \item encuentros entre partículas;
    \item distancias mínimas entre objetos móviles.
\end{itemize}

En dos o tres dimensiones debe utilizarse
prioritariamente el método de componentes rectangulares.

Las ecuaciones clásicas desarrolladas
requieren ejes paralelos y velocidades
suficientemente pequeñas frente a
la velocidad de la luz.

La idea fundamental del capítulo es:

\[
\boxed{
\text{El movimiento se describe
respecto de un sistema de referencia.}
}
\]

Un mismo objeto puede tener diferentes
posiciones, velocidades y trayectorias
según el observador, pero las diferentes
descripciones deben relacionarse
mediante ecuaciones físicas coherentes.

\end{resumen}